\section*{Chapitre \ref{ch:b2:genome-diversification} --- Mutations et diversification des génomes}

\begin{solution}{exo:b2:genome-diversification:1}
GAG~: silencieuse (toujours du glutamate), une transition (A$\to$G).
GTA~: faux-sens (valine), une transversion (A$\to$T). TAA~: non-sens
(stop), une transversion (G$\to$T). Insertion d'une base~: un \omterm{def:b2:genome-diversification:mutation}{décalage
du cadre de lecture}, qui réécrit tous les codons situés après elle.
\end{solution}

\begin{solution}{exo:b2:genome-diversification:2}
$U = 5\times 10^{-10}\times 4.6\times 10^{6} = 2.3\times 10^{-3}$~: une
division sur 430 environ produit une \omterm{def:b2:genome-diversification:mutation}{mutation}. Dans $10^{9}$ cellules~:
$2.3\times 10^{6}$ \omterm{def:b2:genome-diversification:mutation}{mutations} nouvelles par génération.
\end{solution}

\begin{solution}{exo:b2:genome-diversification:3}
\omterm{def:b2:genome-diversification:hgt}{Transformation}~: absorption d'ADN libre issu de cellules lysées~; exige
un receveur compétent~; transfère des fragments chromosomiques (gènes de
capsule chez les pneumocoques). Conjugaison~: contact de cellule à
cellule par un pilus~; exige un \omterm{def:b2:unicellular-diversity:prokaryote}{plasmide} conjugatif chez le donneur~;
transfère des \omterm{def:b2:unicellular-diversity:prokaryote}{plasmides} (gènes de résistance) et, depuis une souche Hfr,
des gènes chromosomiques. \omterm{def:b2:genome-diversification:hgt}{Transduction}~: un bactériophage ayant
empaqueté de l'ADN de l'hôte~; transfère un fragment quelconque
(\omterm{def:b2:genome-diversification:hgt}{transduction} généralisée) ou les gènes voisins du site du prophage
(\omterm{def:b2:genome-diversification:hgt}{transduction} spécialisée).
\end{solution}

\begin{solution}{exo:b2:genome-diversification:4}
\omterm{def:b2:genome-diversification:polyploidy}{Autopolyploïde}~: jeux supplémentaires issus d'une même espèce, par un
gamète diploïde (pomme de terre tétraploïde, nombreux trèfles
autotétraploïdes). \omterm{def:b2:genome-diversification:polyploidy}{Allopolyploïde}~: jeux de deux espèces réunis par
hybridation puis doublés (blé tendre, cotonnier, tabac). L'hybride
diploïde AB ne possède pas deux chromosomes semblables~: aucun ne peut
s'apparier en méiose I, les chromosomes ségrègent au hasard et les
gamètes sont déséquilibrés et non viables.
\end{solution}

\begin{solution}{exo:b2:genome-diversification:5}
$2\times 3.2\times 10^{9}\times 1.2\times 10^{-8} = 77$ \omterm{def:b2:genome-diversification:mutation}{mutations}
nouvelles. Dans les séquences codantes~: $77\times 0.015 = 1.2$~;
changeant un acide aminé~: $1.2\times 0.7 = 0.8$ --- la plupart des
enfants portent environ une substitution d'acide aminé nouvelle.
\end{solution}

\begin{solution}{exo:b2:genome-diversification:6}
Événements~: $m(N - N_0) = 2\times 10^{-8}\times 10^{9} = 20$. Cellules
mutantes~: $mN\ln(N/N_0) = 20\times \ln 10^{6} = 20\times 13.8 = 276$.
$p_0 = e^{-20} \approx 2\times 10^{-9}$~: toutes les cultures ont des
mutants, si bien que compter les cultures qui n'en ont pas ne mesure
rien. Il faut $m(N -
N_0) \approx 1$, c'est-à-dire $N \approx 5\times 10^{7}$, où $p_0 \approx
0.37$, et quelques dizaines de cultures donnent $m$ à un facteur deux
près.
\end{solution}

\begin{solution}{exo:b2:genome-diversification:7}
L'uracile n'a pas sa place dans l'ADN~: une uracile-ADN glycosylase peut
donc retirer sans ambiguïté tout uracile qu'elle rencontre. La thymine
est une base normale~: après la désamination d'une 5-méthylcytosine, le
mésappariement G$\cdot$T est détecté, mais la machinerie de réparation
n'a aucun moyen de savoir quel brin est erroné et, une fois sur deux,
elle «~corrige~» le G, ce qui fixe une transition C$\to$T. Les sites
CpG méthylés mutent donc environ dix fois plus vite que les autres
sites, et au cours de l'évolution ils ont été convertis en TpG et en
CpA~: les génomes de mammifères ne contiennent qu'environ un cinquième
des CpG que prédit leur composition en bases.
\end{solution}

\begin{solution}{exo:b2:genome-diversification:8}
\qty{46}{kb} par minute. Positions~: 368, 782, 1150 et
\qty{1840}{kb}~; distances 414, 368 et \qty{690}{kb}.
\end{solution}

\begin{solution}{exo:b2:genome-diversification:9}
La deuxième copie d'un chromosome s'apparie avec la première copie de
l'autre~; un crossing-over entre elles donne une chromatide portant la
copie~1, la copie~2 et la copie~2$'$ (trois gènes) et une chromatide ne
portant que la copie~1$'$. Une personne ayant reçu un chromosome à un
gène de chacun de ses parents possède deux gènes $\alpha$ au lieu de
quatre ($-\alpha/-\alpha$)~: les chaînes $\alpha$
sont produites à moitié du rythme normal, les hématies sont petites et
pâles, avec une anémie légère ou absente --- c'est le trait
$\alpha$-thalassémique.
\end{solution}

\begin{solution}{exo:b2:genome-diversification:10}
Gamètes~: $n = 21$ (un jeu de chacun de A, B, D). Hybrides~: AB, 14
chromosomes~; ABD, 21. Blé tendre $\times$ amidonnier~: AABBD, 35
chromosomes --- les jeux A et B s'apparient, le jeu D n'a pas de
partenaires, et la plante est largement stérile. Chaque doublement a
donné à chaque chromosome un homologue identique~: des bivalents se
forment et les gamètes sont équilibrés~; un croisement en retour avec un
parent produit des plantes dont un jeu reste non apparié et dont les
gamètes sont déséquilibrés, si bien que l'hexaploïde est isolé
reproductivement de ses ancêtres --- une espèce en une seule étape.
\end{solution}

\begin{solution}{exo:b2:genome-diversification:11}
Croissance logistique~: $R(t) = N/\bigl(1 + (N - 1)e^{-\gamma Nt}\bigr)$.
La moitié de la population est atteinte quand $(N - 1)e^{-\gamma Nt} = 1$~:
$t = \ln(N -
1)/\gamma N = 20.7/0.5 = \qty{41}{h}$. Sans l'antibiotique, les cellules
sans \omterm{def:b2:unicellular-diversity:prokaryote}{plasmide} croissent \qty{5}{\%} plus vite~; le rapport des porteurs
aux non-porteurs décroît donc comme $e^{-st}$ avec $s = 0.05\ln 2 =
\qty{0.035}{h^{-1}}$ (\omterm{def:b2:unicellular-diversity:division}{temps de génération} \qty{1}{h})~: passer de
\qty{99}{\%} à \qty{1}{\%} prend $\ln(99^2)/0.035 = \qty{260}{h}$, onze
jours --- et la poursuite des transferts peut maintenir indéfiniment le
\omterm{def:b2:unicellular-diversity:prokaryote}{plasmide} à basse fréquence, prêt pour la prochaine cure du médicament.
\end{solution}

\begin{solution}{exo:b2:genome-diversification:12}
Le \omterm{thm:b2:genome-diversification:luria}{test de fluctuation} a montré que les \omterm{def:b2:genome-diversification:mutation}{mutations} de résistance
apparaissent avant l'agent sélectif et sans lui, à des moments
aléatoires~: la \omterm{def:b2:genome-diversification:mutation}{mutation} est donc aveugle à ce qui serait utile~;
l'étalement sur répliques a montré la même chose sans jamais exposer
les cellules sélectionnées. L'évolution n'est pas aléatoire, car la
sélection trie ces variants aveugles selon leurs conséquences, de
manière constante et cumulative. Le seul endroit où le hasard est
lui-même façonné est le taux~: les lignées qui commettent trop ou trop
peu d'erreurs perdent, si bien que le \omterm{thm:b2:genome-diversification:rate}{taux de mutation} est un compromis
adapté --- mais ce que fait chaque \omterm{def:b2:genome-diversification:mutation}{mutation} reste imprévisible et non
choisi.
\end{solution}

\begin{solution}{pb:b2:genome-diversification:1}
\textbf{1.} Somme 119, moyenne $5.95$~; moyenne des carrés
$11489/20 = 574$~; variance $574 - 35 = 540$.
\textbf{2.} Somme 324, moyenne $16.2$~; moyenne des carrés
$5462/20 = 273$~; variance $273 - 262 = 11$.
\textbf{3.} Les échantillons d'une même culture suivent une loi de
Poisson (variance proche de la moyenne)~: ce sont des tirages au hasard
dans une même fraction fixe de mutants. Les cultures indépendantes ont
une variance quatre-vingt-dix fois supérieure à la moyenne~: leurs
mutants sont apparus à des moments différents et aléatoires de la
croissance, et plus la \omterm{def:b2:genome-diversification:mutation}{mutation} est précoce, plus son clone est grand.
\textbf{4.} Un jackpot~: une \omterm{def:b2:genome-diversification:mutation}{mutation} survenue au cours de l'une des
toutes premières divisions, dont le clone a ensuite crû avec la culture
jusqu'à une centaine de cellules.
\textbf{5.} Quatorze sur vingt~: $p_0 = 0.70$~; $m(N - N_0) = -\ln 0.70
= 0.36$~; $m = 0.36/2\times 10^{8} = \qty{1.8e-9}{}$ par division.
\textbf{6.} $1.8\times 10^{-9}\times 2\times 10^{8}\times \ln(2\times
10^{6}) = 0.36\times 14.5 = 5.2$ cellules mutantes~; observé~: $5.95$
--- accord compatible avec le bruit de vingt cultures.
\textbf{7.} Aucun échantillon n'était vide (tous partagent les mutants
de la culture mère)~: $p_0 = 0$ ne donne aucune estimation~; les
échantillons mesurent la fraction de mutants d'une seule culture, qui
dépend du moment où ses \omterm{def:b2:genome-diversification:mutation}{mutations} se sont trouvées survenir.
\textbf{8.} $u = m/20 = \qty{9e-11}{}$ par paire de bases et par
réplication.
\textbf{9.} $U = 9\times 10^{-11}\times 4.6\times 10^{6} =
\qty{4.1e-4}{}$.
\textbf{10.} $10^{12}$ cellules $\times$ $4.1\times 10^{-4}$ $= 4.1\times
10^{8}$ \omterm{def:b2:genome-diversification:mutation}{mutations} nouvelles.
\textbf{11.} $3\times 4.6\times 10^{6} = 1.4\times 10^{7}$ changements
possibles~; chacun apparaît $4.1\times 10^{8}/1.4\times 10^{7} \approx 30$
fois.
\textbf{12.} Tout changement d'une seule base, y compris chacun de ceux
qui confèrent la résistance à un médicament mis en échec par une seule
\omterm{def:b2:genome-diversification:mutation}{mutation}, est présent dans des dizaines de cellules avant l'ajout du
médicament~; celui-ci ne fait que les sélectionner.
\textbf{13.} Deux \omterm{def:b2:genome-diversification:mutation}{mutations} indépendantes dans une même cellule
surviennent avec une probabilité $m^2 \approx (2\times 10^{-9})^2 = 4\times 10^{-18}$ par division~: dans
$10^{12}$ cellules, $4\times 10^{-6}$ par génération --- une cellule
doublement résistante ne préexiste pratiquement jamais, et c'est
pourquoi la tuberculose se traite par plusieurs médicaments à la fois.
\textbf{14.} Avec $a = N - 1$ et $E = e^{-\gamma Nt}$~: $R = N/(1 +
aE)$, $\mathrm{d}R/\mathrm{d}t = \gamma N\cdot NaE/(1 + aE)^2$~; et
$\gamma R(N - R) = \gamma\,\frac{N}{1 + aE}\cdot\frac{NaE}{1 + aE}$,
qui est identique~; $R(0) = N/(1 + N - 1) = 1$.
\textbf{15.} $(N - 1)E = 1$~: $t = \ln(N - 1)/\gamma N = 20.7/0.5 =
\qty{41}{h}$.
\textbf{16.} $1 + aE = 1/0.99$, $aE = 0.0101$~: $t = \ln(10^{9}/0.0101)
/0.5 = (20.7 + 4.6)/0.5 = \qty{51}{h}$.
\textbf{17.} $\gamma (N/2)^2 = \gamma N\cdot N/4 = 0.5\times 2.5\times
10^{8} = 1.25\times 10^{8}$ transferts par heure.
\textbf{18.} Taux de croissance $\ln 2 = \qty{0.693}{h^{-1}}$ et
\qty{0.762}{h^{-1}}~; le rapport mutants/sensibles croît comme $e^{st}$
avec $s = \qty{0.069}{h^{-1}}$~; passer de $1/N$ à 1 prend $\ln N/s =
20.7/0.069 = \qty{300}{h}$, douze jours.
\textbf{19.} Le transfert est proportionnel au produit du nombre de
porteurs par celui des non-porteurs et va au rythme des rencontres, non
des divisions~; la descendance est limitée par l'avantage de croissance
du mutant, qui est faible. Un \omterm{def:b2:unicellular-diversity:prokaryote}{plasmide} traverse la population en deux
jours, un mutant favorisé en deux semaines.
\textbf{20.} $U = 4.1\times 10^{-2}$~; \omterm{def:b2:genome-diversification:mutation}{mutations} nuisibles~:
$0.041\times 0.88\times
0.3 = 0.011$ par génération.
\textbf{21.} Moyenne $2.2$~: $e^{-2.2} = 0.11$, un descendant sur neuf
indemne.
\textbf{22.} Type sauvage~: $1.1\times 10^{-4}$ par génération, $0.022$
en 200~: $e^{-0.022} = 0.98$.
\textbf{23.} Moyenne $520$ cellules mutantes~; $p_0 = e^{-36} \approx
10^{-16}$~: aucune culture vide.
\textbf{24.} Sous une sélection intense et changeante (la rotation des
antibiotiques d'un hôpital), l'apport centuplé de variants nouveaux
l'emporte sur le fardeau, et les mutateurs sont entraînés en
auto-stop avec les résistances qu'ils engendrent. Dans un environnement
stable, il n'y a rien à gagner et une \omterm{def:b2:genome-diversification:mutation}{mutation} nuisible par centaine de
générations à perdre~: les lignées mutatrices sont éliminées.
\textbf{25.} $m = \qty{1.8e-9}{}$ par division vers la résistance, $u =
\qty{9e-11}{}$ par paire de bases et par réplication,
$U = \qty{4.1e-4}{}$ \omterm{def:b2:genome-diversification:mutation}{mutation} par génome et par réplication~; le
\omterm{def:b2:unicellular-diversity:prokaryote}{plasmide} atteint la moitié de l'intestin en \qty{41}{h} et \qty{99}{\%}
de celui-ci en \qty{51}{h}.
\end{solution}
